\section{Introduction}
\label{sec:introduction}

A useful personal agent must remember. A calendar correction should survive the conversation in which it was made; a recurring accessibility preference should not need to be repeated; a food allergy mentioned weeks earlier may matter to a shopping or restaurant task. Systems such as MemGPT, MemoryBank, and reflective agents show how persistent stores and generated summaries can extend an agent beyond a finite context window \citep{packer2023memgpt,zhong2024memorybank,shinn2023reflexion}. On a phone, however, ``memory'' is not merely a retrieval-quality component. It is a durable, multi-principal data service positioned among applications, tools, system APIs, cloud models, and one person whose disclosures were made in different contexts.

The security mismatch is concrete. Android AppSearch offers a system-wide index and package/schema visibility controls, and AppFunctions exposes typed callable functionality \citep{android2026appsearch,android2026visibility,android2026appfunctions}. These mechanisms identify a caller or function, but they do not natively express that a fact acquired while arranging health care may be used by a medication-reminder tool but not by an advertising or shopping tool; that the authorization expires on a date; or that a summary derived jointly from health and calendar items must retain both restrictions. Conversely, application-level memory libraries normally expose namespaces and retrieval filters chosen by the agent application. A compromised prompt, an over-sharing model, or a malicious tool therefore operates on the same side of the boundary as the data it is supposed to protect.

Recent evidence makes this boundary timely. MEXTRA demonstrates black-box extraction from agent memories \citep{wang2025mextra}; MINJA injects records through ordinary interaction \citep{dong2025minja}; ADAM adapts extraction queries to the apparent memory distribution \citep{lyu2026adam}; and SPORE persists reactivation payloads and uses a malicious tool channel to export records despite user-level isolation \citep{gao2026spore}. ToolPrivacyBench separately shows why task success is not a privacy metric: a tool-using trajectory can complete correctly while sending unnecessary private atoms to intermediate tools \citep{hu2026toolprivacy}. Memory contamination is durable as well: conversational poisoning and multi-record composition can steer later behavior \citep{wang2026hijacking,gao2026mempoison,xie2026memevobench}. These works do not imply that every memory deployment is exploitable in the same way, but they invalidate the assumption that a user namespace is a sufficient authorization boundary.

\subsection{A purpose-bound service boundary}

This paper develops \ms, an Android service that mediates persistent personal-agent memory outside the model. A memory document carries a label containing its provenance, allowed purposes, allowed app/tool consumers, sensitivity class, creation time, and expiry. An agent never selects an arbitrary purpose string and never queries the backing store directly. Instead, a system component issues a signed retrieval token from the active task and AppFunction execution context. The retrieval monitor first constructs an authorized view and only then invokes lexical or vector ranking. Thus an out-of-scope item is not merely down-ranked or hidden by an instruction; its plaintext is absent from the candidate set visible to the ranker and response model.

Derived memories create the less familiar systems problem. An LLM may compress ten messages into a summary, infer a stable preference, or write a reflection used by later tasks. If a new item combines sources, assigning it the least restrictive source label would launder information across contexts. MemScope instead computes the meet of source authorizations: set-valued permissions are intersected, sensitivity is joined upward, and expiry becomes no later than the earliest source expiry. The service requires cited source identifiers and verifies structural attribution before accepting a caller-proposed derivation. A semantic checker may reject unsupported content, but it is not trusted to widen policy. When attribution is uncertain, admission fails closed or the derived item receives a quarantine label.

Lifetime and resource pressure are also authorization concerns. Deleting an index entry in a later maintenance pass is insufficient if an expired item remains readable between expiry and compaction. MemScope checks a trusted monotonic time on every retrieval and performs background deletion as space reclamation rather than as the primary enforcement mechanism. Under storage pressure, policy-aware eviction preserves pinned safety-critical facts when policy permits, while never extending their lifetime or consumer scope. This differs from treating an LLM context cache as the durable record: recent work shows that safety-critical facts can disappear when compression is evaluated without the future question \citep{zhu2026allergy}.

\subsection{Scope and claims}

The formal claims are intentionally narrower than ``memory privacy.'' Given an uncompromised OS service, authentic tokens, complete mediation, correct principal binding, and cryptographic compartment assumptions stated in \cref{sec:threat}, we show that any plaintext item returned by the MemScope API satisfies the declared authorization predicate. We show that a derived label cannot authorize a purpose, consumer, app, tool, or time beyond all its sources. We show that an item whose lifetime has elapsed is not returned, even before physical compaction. These are safety properties of the service interface. They do not prevent an authorized consumer from memorizing or re-transmitting data, do not hide access patterns, and do not provide semantic non-inference over a sequence of individually authorized responses. DP-MemView addresses that complementary cumulative-inference problem under a different interface contract \citep{kim2026dpmemview,dwork2014algorithmic}.

No completed MemScope implementation or measurement dataset was present in the source project. Accordingly, this manuscript does not claim an achieved extraction rate, latency, energy cost, accuracy, or statistical improvement. It contributes a reproducible system design, formal core, prototype plan, and experimental protocol. Conceptual plots make hypotheses inspectable but are visibly labeled as non-results.

\subsection{Contributions}

The work makes five contributions directly expanding the source concept:

\begin{enumerate}[leftmargin=*]
  \item \textbf{Labeled mobile-agent memory.} We define a persistent document abstraction whose enforceable dimensions are provenance, purpose, app/tool consumer, sensitivity, and lifetime, coupled to a retrieval-token protocol rooted in Android task and AppFunction context.
  \item \textbf{Conservative derivation.} We specify source-cited writes, meet-based label propagation, attribution validation, and fail-closed quarantine for summaries and reflections. This prevents the policy widening that would otherwise accompany LLM-based compression.
  \item \textbf{Lifecycle enforcement under mobile constraints.} We separate retrieval-time expiry from physical deletion and introduce a policy-aware eviction objective that preserves utility and designated safety value without changing authorization.
  \item \textbf{A precise security argument.} Under stated assumptions, we prove admission soundness, derived-label monotonicity, interface-level cross-scope non-disclosure, expiry safety, and explicit complexity bounds. We also state counterexamples outside those assumptions.
  \item \textbf{An on-device evaluation protocol.} We define matched experiments on AOSP phones against unscoped memory, identity namespaces, AppSearch visibility, prompt-only instructions, and GAAP-style IFC, using memory extraction/poisoning attacks and long-horizon utility workloads. Every presently unmeasured value remains a placeholder.
\end{enumerate}

\subsection{Paper organization}

\Cref{sec:related} positions MemScope against memory systems, attacks, information-flow control, and mobile platforms. \Cref{sec:preliminaries,sec:system,sec:threat,sec:problem} define the objects and adversary. \Cref{sec:method,sec:algorithm} give the architecture and algorithms. \Cref{sec:theory} presents the formal results; \cref{sec:implementation} maps them to AOSP. \Cref{sec:evaluation} specifies the experiments. \Cref{sec:discussion,sec:limitations} discuss deployment and limits before \cref{sec:conclusion} concludes.
